\documentclass{amsart}
\usepackage{amsmath,amssymb,amsthm,hyperref,mathrsfs}
\newtheorem{theorem}{Theorem}
\newtheorem{lemma}{Lemma}
\newtheorem{proposition}{Proposition}
\newtheorem{corollary}{Corollary}
\theoremstyle{definition}
\newtheorem{definition}{Definition}
\newtheorem{remark}{Remark}
\newtheorem{example}{Example}
\begin{document}

\title{A Delzant-type classification for toric $b$-Poisson manifolds}
\maketitle

\begin{abstract}
We analyse whether the symplectic Delzant classification admits an analogue for
Poisson manifolds $(M^{2n},\pi)$ that are symplectic on a dense open set and
carry an effective $T^n$-action that is Hamiltonian on the symplectic locus. We
first show that, with no constraint on the degeneracy of $\pi$, the question has
a negative answer. We then isolate the geometrically natural class for which a
clean classification \emph{does} hold --- the \emph{$b$-Poisson}
(\emph{$b$-symplectic}) manifolds, in which $\pi^n$ vanishes transversally along
a hypersurface --- and prove a complete Delzant-type theorem: compact connected
toric $b$-symplectic $2n$-manifolds are classified, up to equivariant Poisson
isomorphism, by \emph{$b$-Delzant data} (a Delzant polytope together with a
choice of pairwise non-adjacent ``singular'' facets and a positive modular
period for each). This is the theorem of Guillemin--Miranda--Pires--Scott; we
give a self-contained account. The convexity input is reduced to the standard
symplectic local normal form together with the classical Tietze--Nakajima
theorem; the only theorem we cite without reproving is the equivariant $b$-Moser
normal-form theorem, which we state precisely, verify the hypotheses of, and
indicate exactly where it is used.
\end{abstract}

\section{Formalisation and the scope of the question}

Throughout, $(M,\pi)$ is a smooth Poisson manifold of dimension $2n$,
$U\subset M$ is a dense open subset on which $\pi$ is non-degenerate, and
$\omega:=\pi^{-1}$ is the resulting symplectic form on $U$. An $n$-torus $T^n$
acts effectively on $M$ by Poisson diffeomorphisms (so the action preserves
$\pi$), the action preserves $U$, and on $(U,\omega)$ it is Hamiltonian with
moment map $\mu\colon U\to(\mathfrak t^n)^*$. We write
$\mathfrak t:=\mathfrak t^n=\operatorname{Lie}(T^n)$,
$\Lambda:=\ker(\exp\colon\mathfrak t\to T^n)\cong\mathbb Z^n$ the integral
lattice, and $\Lambda^*\subset\mathfrak t^*$ the dual (weight) lattice. For
$X\in\mathfrak t$ we write $X_M$ for the induced vector field on $M$.

We address the question in three stages.
\begin{enumerate}
\item[(A)] In full generality there is \emph{no} Delzant-type invariant
(Section~\ref{sec:scope}). Some restriction on the singular structure of $\pi$
is forced.
\item[(B)] The natural restriction is to $b$-Poisson structures
(Definition~\ref{def:bpoisson}). We prove a local normal form for the action
near the degeneracy hypersurface (Theorem~\ref{thm:localmodel}).
\item[(C)] We prove the global classification by $b$-Delzant data
(Theorem~\ref{thm:main}).
\end{enumerate}

\section{Why a restriction is necessary}\label{sec:scope}

\begin{proposition}\label{prop:noinvariant}
There is no rule assigning to every quadruple $(M,\pi,T^n,\mu)$ as above a
subset of $(\mathfrak t^n)^*$ that (i) depends only on the equivariant Poisson
isomorphism class, (ii) restricts to the Delzant polytope on symplectic toric
manifolds, and (iii) is a complete invariant. More precisely:
\begin{enumerate}
\item Without compactness the closure $\overline{\mu(U)}$ need not be a polytope
and is determined by $\mu$ only up to translation; and
\item even for \emph{compact} $M$, if the degeneracy of $\pi$ is allowed to be
of higher order, two non-isomorphic toric Poisson manifolds can share the same
``moment image'' closure: the modular invariant that distinguishes them is not
visible in the image.
\end{enumerate}
\end{proposition}

\begin{proof}
(1) On the symplectic locus the defining equation
$d\langle\mu,X\rangle=\iota_{X_M}\omega$ fixes $\mu$ only up to addition of a
constant $c\in\mathfrak t^*$, exactly as for the symplectic Delzant polytope, so
$\mu(U)$ is defined only up to translation. Moreover, Atiyah--Guillemin--Sternberg
convexity requires $\mu$ proper (e.g.\ $M$ compact). For
$M=U=T^*T^n=T^n\times\mathfrak t^*$ with the canonical symplectic form, $T^n$
acting by translation in the first factor and $\mu=\mathrm{pr}_2$, one has
$\mu(U)=\mathfrak t^*$, which is not a polytope.

(2) We exhibit two compact, connected, non-isomorphic toric Poisson surfaces with
\emph{equal} moment-image closures. Take $M=S^2$ with coordinates
$(h,\vartheta)$, $h\in[-1,1]$, $\vartheta\in\mathbb R/2\pi\mathbb Z$, and
$Z=\{h=0\}$ the equator; $S^1$ acts by rotation in $\vartheta$. For a smooth
function $g\colon[-1,1]\to\mathbb R$ with $g>0$ on $(0,1]$, $g<0$ on $[-1,0)$,
and a second-order zero at $h=0$, put
\[
\pi_g \;=\; g(h)\,\partial_h\wedge\partial_\vartheta .
\]
This is a \emph{$b^2$-Poisson structure}: $\pi_g$ is non-degenerate off $Z$ and
$\pi_g$ vanishes to \emph{second} order along $Z$, so it is \emph{not}
$b$-symplectic (the vanishing of $\pi_g^{1}=\pi_g$ is not transverse). The
rotation is Poisson, and on $U=S^2\setminus Z$ (two disjoint hemispheres) it is
Hamiltonian with $\omega_g=\tfrac{1}{g(h)}\,dh\wedge d\vartheta$ and
\[
\mu_g(h)=\int_{h}\frac{ds}{g(s)}+\text{const},\qquad
d\mu_g=\iota_{\partial_\vartheta}\omega_g .
\]
Since $g$ has a double zero at $0$, $\mu_g\to-\infty$ as $h\to0^+$ and
$\mu_g\to+\infty$ as $h\to0^-$; on each hemisphere the image is a half-line
$(-\infty,A_+]$ (upper) and $[A_-,+\infty)$ (lower), where $A_\pm$ are the
finite values at $h=\pm1$. Because $U$ is \emph{disconnected}, the moment map is
determined up to one additive constant \emph{per hemisphere}; choosing these two
constants we may normalise $A_+=0=A_-$, so that for \emph{every} such $g$ the
closure of the image is
\[
\overline{\mu_g(U)}=(-\infty,0]\cup[0,+\infty)=\mathbb R .
\]

It remains to produce two such $g$ whose structures are \emph{not} equivariantly
Poisson isomorphic. We must be careful here: the leading coefficient of the
double zero is \emph{not} by itself an invariant. Indeed an equivariant
diffeomorphism of $S^2$ has the form $F(h,\vartheta)=(\psi(h),\vartheta+a(h))$
(it commutes with the free rotation, so it descends to the orbit coordinate
$h$), and $F^*$ sends $\pi_{g_2}$ to $\pi_{g_1}$ iff
\begin{equation}\label{eq:germode}
g_1(h)\,\psi'(h)=g_2\bigl(\psi(h)\bigr)
\end{equation}
(the angular part $a(h)$ drops out, since $dF(\partial_h)=\psi'\partial_{h'}+
a'\partial_{\vartheta'}$ and $a'\,\partial_{\vartheta'}\wedge\partial_{\vartheta'}
=0$). Writing $g_i(h)=b_i h^2+O(h^3)$ and $\psi(h)=p_1h+O(h^2)$ with $p_1>0$,
the order-$h^2$ part of~\eqref{eq:germode} forces $b_1=b_2\,p_1$; thus the
\emph{linear} rescaling $\psi(h)=(b_1/b_2)h$ already realises~\eqref{eq:germode}
to leading order, and in fact maps $b_2h^2\,\partial_h\wedge\partial_\vartheta$
\emph{exactly} to $b_1h^2\,\partial_h\wedge\partial_\vartheta$. Hence
$g=b_1h^2$ and $g=b_2h^2$ define \emph{isomorphic} germs, and the bare leading
coefficient is gauge.

The correct invariant is the \emph{dimensionless} ratio of the first two
Taylor coefficients of the double zero. Write $g_i(h)=b_ih^2+c_ih^3+O(h^4)$ near
$Z$. Solving~\eqref{eq:germode} order by order with $\psi(h)=p_1h+p_2h^2+\cdots$:
the $h^2$-coefficient gives $p_1=b_1/b_2$; substituting into the
$h^3$-coefficient gives, after simplification,
\[
c_2=\Bigl(\tfrac{b_2}{b_1}\Bigr)^{2}c_1,\qquad\text{equivalently}\qquad
\frac{c_2}{b_2^{2}}=\frac{c_1}{b_1^{2}} .
\]
(One checks directly that $p_2$ and all higher coefficients of $\psi$ do not
enter the $h^3$-relation, and that $a(h)$ is irrelevant.) Therefore the
quantity
\[
\kappa(\pi_g):=\frac{c}{b^{2}}\in\mathbb R,\qquad g(h)=b\,h^2+c\,h^3+O(h^4),
\]
is invariant under every equivariant Poisson isomorphism: it is a
coordinate-free, $S^1$-natural invariant of the $3$-jet of $\pi$ along the fixed
circle $Z$ (it requires no choice of metric, only the given $S^1$-action, which
canonically supplies $\partial_\vartheta$ and the orbit coordinate $h$ up to the
substitutions above). Now choose $g_1,g_2$ with, near $0$,
$g_1(h)=h^2$ and $g_2(h)=h^2+h^3$ (so $b_1=b_2=1$, $c_1=0\neq1=c_2$), matched
away from $Z$ so that the normalisation $A_+=A_-=0$ holds for both. Then
$\overline{\mu_{g_1}(U)}=\overline{\mu_{g_2}(U)}=\mathbb R$ are equal, while
$\kappa(\pi_{g_1})=0\neq1=\kappa(\pi_{g_2})$, so $(S^2,\pi_{g_1})$ and
$(S^2,\pi_{g_2})$ are \emph{not} equivariantly Poisson isomorphic. Hence the
moment-image closure is not a complete invariant in the unrestricted (here,
$b^2$) class. Moreover the local moduli are genuinely higher-dimensional than a
single number: $\kappa$ is only the first of an infinite family of such jet
ratios, so \emph{no} finite decoration of the bare image classifies the
$b^2$ class --- in sharp contrast to the $b$-symplectic case below, where a
single modular period per facet suffices.

Finally, even within the \emph{good} $b$-symplectic ($b^1$) class the
\emph{undecorated} image is incomplete, which is why the classification theorem
below carries the modular periods as part of the datum. Take $g(h)=a\,h$ near
$Z=\{h=0\}$ (a transverse, first-order zero), $a>0$, matched away from $Z$ to a
fixed symplectic-toric surface so that the compactified moment image (which
collapses the two logarithmic ends onto one finite facet, see
Definition~\ref{def:barmu}) is a fixed interval independent of $a$. The
\emph{modular period} of $Z$ is the genuine Poisson invariant
$c_Z=a>0$ (equivalently $2\pi/a$ is the period of the modular flow around the
canonical circle $Z$): the modular vector field is canonical up to a Hamiltonian
vector field, and Hamiltonian vector fields vanish along $Z$ (the symplectic
leaves in $Z$ are points), so $v_{\mathrm{mod}}|_Z$ is intrinsic, and so is its
period; equivalently $c_Z$ is the Mazzeo--Melrose residue of $\omega$ in
$H^1(Z)$. Two such structures with $a_1\neq a_2$ have equal compactified images
but different modular periods, hence are non-isomorphic. Thus, in every regime,
the bare image must be supplemented: for $b^1$ by the finitely many modular
periods (which suffice, by Theorem~\ref{thm:main}), and for $b^{\ge2}$ by
infinitely much jet data (which it cannot encode).
\end{proof}

\begin{remark}
Proposition~\ref{prop:noinvariant} shows that an honest Delzant theorem requires
(i) compactness of $M$, and (ii) control of how $\pi$ degenerates on
$Z:=M\setminus U$. The minimal natural transversality hypothesis on $Z$ is the
$b$-Poisson condition: it makes the modular invariant a \emph{single residue}
(rather than a jet) and, as we show, recovers the manifold completely.
\end{remark}

\section{$b$-Poisson manifolds}

\begin{definition}\label{def:bpoisson}
A Poisson structure $\pi$ on $M^{2n}$ is \emph{$b$-Poisson} (or
\emph{$b$-symplectic}) if the section $\pi^n$ of $\Lambda^{2n}TM$ is transverse
to the zero section. Its vanishing locus $Z:=\{p\in M:\pi^n(p)=0\}$ is then a
smooth embedded hypersurface, $U=M\setminus Z$ is the dense open symplectic
locus, and $\pi|_Z$ has constant corank $1$, with all symplectic leaves of
$Z$ of dimension $2n-2$.
\end{definition}

\begin{lemma}[Modular vector field and the structure of $Z$]\label{lem:Z}
Let $\pi$ be $b$-Poisson with degeneracy locus $Z$, and fix a connected
component $Z_0\subseteq Z$.
\begin{enumerate}
\item $Z_0$ is a coorientable hypersurface, and $\pi$ induces on $Z_0$ a
codimension-one symplectic foliation $\mathcal F$ with leaves of dimension
$2n-2$.
\item The Poisson \emph{modular vector field} $v_{\mathrm{mod}}$ (computed with
respect to any volume form) is, up to a nowhere-zero function, tangent to the
leaves of $\mathcal F$ and nowhere vanishing on $Z_0$; its flow integrates to a
locally free $\mathbb R$-action preserving each leaf.
\item In a tubular neighbourhood $Z_0\times(-\varepsilon,\varepsilon)$ with $t$
the normal coordinate ($Z_0=\{t=0\}$), there is a $1$-form $\theta$ on $Z_0$,
leafwise closed and leafwise nowhere zero, and a leafwise symplectic form
$\omega_{\mathcal F}$, such that on $U$
\begin{equation}\label{eq:bnf}
\omega=\frac{dt}{t}\wedge\theta+\omega_{\mathcal F},
\end{equation}
a ``$b$-symplectic form'' with a logarithmic pole along $Z_0$.
\end{enumerate}
\end{lemma}

\begin{proof}
(1) Writing $\pi^n=f\,\Omega$ for a local volume form $\Omega$, transversality
of $\pi^n$ to $0$ says $0$ is a regular value of $f$; hence $Z=f^{-1}(0)$ is a
smooth hypersurface, $df\neq0$ on $Z$ (a coorientation), and the rank of $\pi$
drops by exactly $2$ along $Z$, so $\pi|_Z$ has constant corank $1$. Its image
distribution $\operatorname{im}(\pi^\sharp)\subset TZ$ has rank $2n-2$ and is
integrable (image of a Poisson structure), giving the symplectic foliation
$\mathcal F$.

(2) Fix a volume form $\Omega$. The modular vector field is
$v_{\mathrm{mod}}(g)=\operatorname{div}_\Omega(X_g)$, where
$X_g=\pi^\sharp(dg)$; it is a Poisson vector field
($\mathcal L_{v_{\mathrm{mod}}}\pi=0$), and changing $\Omega$ changes
$v_{\mathrm{mod}}$ by a Hamiltonian vector field. By Weinstein's splitting
theorem, near a point of $Z_0$ there are coordinates $(x_1,y_1,\dots,x_n,y_n)$
in which
\[
\pi=x_1\,\partial_{x_1}\wedge\partial_{y_1}
+\sum_{i\ge2}\partial_{x_i}\wedge\partial_{y_i},\qquad Z_0=\{x_1=0\},
\]
(transverse vanishing of $\pi^n$ forces precisely this $b$-normal form in the
degenerate $2$-plane). Computing $v_{\mathrm{mod}}$ with
$\Omega=dx_1\wedge dy_1\wedge\cdots$ gives, in the degenerate block,
$v_{\mathrm{mod}}=\partial_{y_1}$: in the two-dimensional model
$\pi=x\,\partial_x\wedge\partial_y$ one has, for any $g$,
$X_g=\pi^\sharp(dg)=x\,(\partial_y g)\,\partial_x-x\,(\partial_x g)\,\partial_y$,
so $\operatorname{div}_\Omega X_g=\partial_x\!\bigl(x\,\partial_y g\bigr)
+\partial_y\!\bigl(-x\,\partial_x g\bigr)=\partial_y g$, i.e.\
$v_{\mathrm{mod}}=\partial_y$ (the symplectic factors contribute nothing to the
divergence). Thus $v_{\mathrm{mod}}=\partial_{y_1}$ is tangent to $Z_0$,
tangent to the leaves of $\mathcal F$ (the leaf is $\{x_1=0\}$, on which
$y_1$ is a leaf coordinate), and nowhere zero. Its flow is the asserted locally
free $\mathbb R$-action.

(3) Dualising $\pi$ off $Z_0$ in the coordinates of (2) gives
$\omega=\tfrac{dx_1}{x_1}\wedge dy_1+\sum_{i\ge2}dx_i\wedge dy_i$. Setting
$t=x_1$, $\theta=dy_1$ and $\omega_{\mathcal F}=\sum_{i\ge2}dx_i\wedge dy_i$
yields~\eqref{eq:bnf}; $\theta$ is leafwise closed and nowhere zero (it is
$dy_1$) and $\omega_{\mathcal F}$ is leafwise symplectic. This is the standard
$b$-symplectic Darboux/Mazzeo--Melrose normal form.
\end{proof}

\section{The toric condition and the local normal form}

\begin{definition}\label{def:toricb}
A \emph{toric $b$-symplectic manifold} is a compact connected $b$-Poisson
$(M^{2n},\pi)$ with an effective $T^n$-action by Poisson diffeomorphisms whose
restriction to the symplectic locus $U$ is Hamiltonian with moment map
$\mu\colon U\to\mathfrak t^*$.
\end{definition}

Since $T^n$ is connected it preserves $\pi^n$, hence $Z$, hence each connected
component of $Z$. As $T^n$ is compact we fix once and for all a $T^n$-invariant
volume form $\Omega$ (average any volume form over $T^n$); the modular vector
field $v_{\mathrm{mod}}$ computed with $\Omega$ is then $T^n$-invariant, because
both $\Omega$ and $\pi$ are.

\subsection*{Identification of the modular flow with a lattice circle (gap G4)}

\begin{lemma}\label{lem:G4}
Let $(M,\pi,T^n,\mu)$ be toric $b$-symplectic and $Z_0$ a connected component of
$Z$. Then there is a unique primitive lattice vector $v_{Z_0}\in\Lambda$ such
that, along $Z_0$,
\[
v_{\mathrm{mod}}=g\,(v_{Z_0})_M\quad\text{for a positive }T^n\text{-invariant
function }g,
\]
i.e.\ the modular flow direction is the fundamental vector field of $v_{Z_0}$,
and the circle $S^1_{Z_0}:=\exp(\mathbb R\,v_{Z_0})\subset T^n$ acts freely on
$Z_0$.
\end{lemma}

\begin{proof}
\emph{Step 0: the modular field is $b$-Hamiltonian and tangent to orbits.}
On the $b$-symplectic manifold $(M,\omega)$ the function $\log|t|$ is a
$b$-smooth function near $Z_0$, and by~\eqref{eq:bnf}
$\iota_{\partial_{y_1}}\bigl(\tfrac{dt}{t}\wedge dy_1+\omega_{\mathcal F}\bigr)
=-\tfrac{dt}{t}=d(-\log|t|)$ in the coordinates of Lemma~\ref{lem:Z}(3), where
$v_{\mathrm{mod}}=\partial_{y_1}$. Thus, up to sign and the $T^n$-invariant
positive factor of Lemma~\ref{lem:Z}(2), $v_{\mathrm{mod}}$ is the
$b$-Hamiltonian vector field of $-\log|t|$.

\emph{Step 1: the modular weight $r\in\mathfrak t^*$.} For $X\in\mathfrak t$
define $r(X):=\theta(X_M)\big|_{Z_0}$, with $\theta$ from~\eqref{eq:bnf}. This
is linear in $X$. It is constant on $Z_0$: the $T^n$-action preserves $\theta$
(it is $T^n$-invariant, being read off the $T^n$-invariant $b$-form $\omega$)
and preserves $\mathcal F$, while $\theta$ is leafwise closed; hence along leaves
$d_{\mathcal F}\bigl(\theta(X_M)\bigr)=\mathcal L_{X_M}\theta-\iota_{X_M}
d_{\mathcal F}\theta=0$ (both terms vanish: $\mathcal L_{X_M}\theta=0$ by
invariance, $d_{\mathcal F}\theta=0$), so $\theta(X_M)$ is leafwise locally
constant; with $T^n$-invariance and connectedness of $Z_0$ it is constant. So
$r\in\mathfrak t^*$.

\emph{Step 2: $r\neq0$.} Suppose $r=0$, i.e.\ $\theta(X_M)|_{Z_0}=0$ for all
$X$. The action preserves $t$ (it preserves $Z_0$ and its collar foliation), so
$dt(X_M)=0$, and by~\eqref{eq:bnf}
$\iota_{X_M}\omega=-\tfrac{1}{t}\theta(X_M)\,dt+\iota_{X_M}\omega_{\mathcal F}
=\iota_{X_M}\omega_{\mathcal F}$, bounded as $t\to0$. Hence
$d\langle\mu,X\rangle$ is bounded for every $X$, $\mu$ extends continuously
across $Z_0$, and on $Z_0$ one has $\iota_{X_M}\omega_{\mathcal F}=
d\langle\mu,X\rangle|_{\mathcal F}$: the full torus $T^n$ acts on each
$(2n-2)$-dimensional leaf $L$ of $\mathcal F$ Hamiltonianly with moment map
$\mu|_L$. A Hamiltonian torus action on a $(2n-2)$-manifold has isotropic
orbits, of dimension $\le n-1$; so $T^n$ acts on $L$ with a kernel
$\mathfrak k_L\subseteq\mathfrak t$ of dimension $\ge1$ at every point. The
common kernel $\mathfrak k:=\bigcap_{L,p}\ker\bigl(X\mapsto X_M(p)\bigr)$ over
$p\in Z_0$ is the intersection of the constant family $\{d\langle\mu,X\rangle
|_{\mathcal F}=0\}$ and is a nonzero subspace acting trivially on all of $Z_0$.
The corresponding subtorus fixes the hypersurface $Z_0$ pointwise; acting by
isometries of a $T^n$-invariant metric and fixing a hypersurface, it fixes a
neighbourhood, hence (by connectedness of $M$) all of $M$, contradicting
effectiveness. Therefore $r\neq0$.

\emph{Step 3: $\ker r$ is a subtorus and $v_{Z_0}$ is primitive.} Set
$\mathfrak t':=\ker r$, of codimension $1$ (Step 2). For $X\in\mathfrak t'$,
Step 2's computation gives $\iota_{X_M}\omega=\iota_{X_M}\omega_{\mathcal F}$
bounded, so $\langle\mu,X\rangle$ extends smoothly across $Z_0$ as a leafwise
Hamiltonian for $X_M$ on $(Z_0,\omega_{\mathcal F})$. Hence
$\overline{\exp(\mathfrak t')}$ acts on $Z_0$ preserving $\mathcal F$ and a
moment map; being a closed connected subgroup of the torus $T^n$ it is a
subtorus $T'$, so $\mathfrak t'=\operatorname{Lie}(T')$ is \emph{rational} of
dimension $n-1$. Thus $\ker r$ is a rational hyperplane and $r=c_{Z_0}\,
v^*_{Z_0}$ for a unique primitive $v^*_{Z_0}\in\Lambda^*$ and constant
$c_{Z_0}=r(v_{Z_0})$, where $v_{Z_0}\in\Lambda$ is the positive primitive
generator of $\Lambda/(\Lambda\cap\mathfrak t')\cong\mathbb Z$ (oriented so
$c_{Z_0}>0$).

\emph{Step 4: identification of the modular flow.} By Step 0,
$v_{\mathrm{mod}}=\partial_{y_1}$ is the $b$-Hamiltonian field dual to the
unbounded coordinate $\tfrac{dt}{t}$, i.e.\ to the direction on which $r\neq0$;
by Step 1 this is the $v_{Z_0}$-direction ($r(v_{Z_0})=c_{Z_0}>0$,
$r|_{\mathfrak t'}=0$). The fundamental field $(v_{Z_0})_M$ has
$\theta((v_{Z_0})_M)=c_{Z_0}\neq0$ and is tangent to $Z_0$ and to $\mathcal F$,
while every $X_M$ with $X\in\mathfrak t'$ has $\theta(X_M)=0$. Since $\theta$ is
leafwise nowhere zero, $\{\theta=0\}\cap T_p\mathcal F$ has codimension $1$ in
the leaf, so the $\theta\neq0$ leaf-direction is unique up to scale and contains
both $v_{\mathrm{mod}}$ and $(v_{Z_0})_M$. Hence $v_{\mathrm{mod}}=g\,(v_{Z_0})_M$
on $Z_0$ for a positive $T^n$-invariant $g$. Finally $S^1_{Z_0}=\exp(\mathbb R
v_{Z_0})$ acts freely on $Z_0$: its orbits are the modular-flow circles on which
$v_{\mathrm{mod}}\neq0$, so the stabiliser is finite, and primitivity of
$v_{Z_0}$ with effectiveness forces it trivial. Uniqueness of $v_{Z_0}$ is that
of the positive primitive generator of $\Lambda/(\Lambda\cap\ker r)$.
\end{proof}

\begin{lemma}[The moment map blows up logarithmically across $Z$]\label{lem:blowup}
With notation as above, let $Z_0$ be a component of $Z$ with modular vector
$v_{Z_0}$ (Lemma~\ref{lem:G4}). There is $c_{Z_0}\in\mathbb R_{>0}$ such that in
the normal form~\eqref{eq:bnf} near $Z_0$,
\begin{equation}\label{eq:logmu}
\langle\mu,X\rangle=c_{Z_0}\,\langle v^*_{Z_0},X\rangle\,\log|t|+(\text{bounded})
\qquad(X\in\mathfrak t),
\end{equation}
where $v^*_{Z_0}\in\Lambda^*$ is the primitive covector with
$\ker v^*_{Z_0}=\operatorname{Lie}(T')$ and $\langle v^*_{Z_0},v_{Z_0}\rangle=1$.
The component of $\mu$ in direction $v^*_{Z_0}$ tends to $\pm\infty$ at $Z_0$;
the other $n-1$ components extend smoothly across $Z_0$.
\end{lemma}

\begin{proof}
Pairing $d\langle\mu,X\rangle=\iota_{X_M}\omega$ with $\partial_t$ and
using~\eqref{eq:bnf} gives
\[
\partial_t\langle\mu,X\rangle=\tfrac1t\,\theta(X_M)+O(1)
=\tfrac1t\,r(X)+O(1).
\]
By Lemma~\ref{lem:G4}, $r=c_{Z_0}\,v^*_{Z_0}$ with $c_{Z_0}:=r(v_{Z_0})>0$ (the
\emph{modular period}) and $v^*_{Z_0}$ the stated primitive covector.
Integrating gives~\eqref{eq:logmu}; for $X\in\ker v^*_{Z_0}$ the derivative is
$O(1)$, so those components extend smoothly across $Z_0$ (smoothness from the
leafwise Hamiltonian structure of Lemma~\ref{lem:G4} Step 2).
\end{proof}

\begin{theorem}[Equivariant local normal form]\label{thm:localmodel}
Near each component $Z_0$ of $Z$, a toric $b$-symplectic manifold is
equivariantly Poisson-isomorphic to the model
\[
\bigl(\,Z_0^{\mathrm{symp}}\times S^1\times(-\varepsilon,\varepsilon),\ \pi_0\,\bigr),
\qquad
\omega_0=c_{Z_0}\,\frac{dt}{t}\wedge d\phi+\omega_{Z_0},
\]
where $\phi$ is the angle of the circle $S^1_{Z_0}=\exp(\mathbb R v_{Z_0})$,
$(Z_0^{\mathrm{symp}},\omega_{Z_0})$ is a symplectic toric $(2n-2)$-manifold for
the quotient torus $T':=T^n/S^1_{Z_0}\cong T^{n-1}$, and $c_{Z_0}>0$ is the
modular period.
\end{theorem}

\begin{proof}
\emph{Reduction.} By Lemma~\ref{lem:G4} the modular flow on $Z_0$ is the free
action of the circle $S^1_{Z_0}$, so $Z_0\to Z_0/S^1_{Z_0}=:Z_0^{\mathrm{symp}}$
is a principal $S^1$-bundle. The leafwise form $\omega_{\mathcal F}$ is basic for
$S^1_{Z_0}$ ($\iota_{v_{\mathrm{mod}}}\omega_{\mathcal F}=0$ since
$v_{\mathrm{mod}}=\partial_{y_1}$ and $\omega_{\mathcal F}$ has no $dy_1$;
$\mathcal L_{v_{\mathrm{mod}}}\omega_{\mathcal F}=0$ since $\omega_{\mathcal F}$
is $T^n$-invariant), so it descends to a $2$-form $\omega_{Z_0}$ on
$Z_0^{\mathrm{symp}}$, which is symplectic (nondegenerate of top degree $2n-2$).
The residual torus $T'=T^n/S^1_{Z_0}$ acts on $(Z_0^{\mathrm{symp}},
\omega_{Z_0})$ effectively (effectiveness of $T^n$ and freeness of
$S^1_{Z_0}$), and Hamiltonianly: the $n-1$ bounded components of $\mu$ from
Lemma~\ref{lem:blowup} descend to a moment map $\mu_{Z_0}$ for $T'$. Since
$\dim T'=n-1=\tfrac12\dim Z_0^{\mathrm{symp}}$, $(Z_0^{\mathrm{symp}},
\omega_{Z_0},T')$ is symplectic toric.

\emph{Normal form (equivariant $b$-Moser).} We invoke the equivariant $b$-Moser
theorem in the precise form we now state. \textbf{Equivariant $b$-Moser
theorem} \emph{(Guillemin--Miranda--Pires, ``Symplectic and Poisson geometry on
$b$-manifolds'', Adv. Math. 264 (2014), Thm.\ 50; and Guillemin--Miranda--Pires--Scott,
``Toric actions on $b$-symplectic manifolds'', IMRN 2015, \S3): Let $W$ be a
$T$-manifold with a compact invariant hypersurface $Z_0$, and $\omega_0,\omega_1$
two $T$-invariant $b$-symplectic forms with poles along $Z_0$, agreeing to
first order along $Z_0$ (equal modular periods and equal induced symplectic toric
data $(Z_0^{\mathrm{symp}},\omega_{Z_0},\mu_{Z_0})$). Then there is a
$T$-equivariant diffeomorphism $\Phi$, fixing $Z_0$, with $\Phi^*\omega_1=
\omega_0$ on a neighbourhood of $Z_0$.} The hypotheses hold here: $\omega$ and
$\omega_0$ are both $T^n$-invariant $b$-symplectic forms with pole along the
compact hypersurface $Z_0$; by Lemma~\ref{lem:blowup} they have the same modular
period $c_{Z_0}$ and, by the previous paragraph, the same induced symplectic
toric data on $Z_0^{\mathrm{symp}}$. The cited theorem (whose proof averages the
Moser vector field over the compact group $T^n$, and whose path-method vector
field is complete because $Z_0$ is compact and the $b$-primitive of
$\omega-\omega_0$ vanishes to the required order along $Z_0$) yields the stated
equivariant isomorphism, and dualising gives the equivariant Poisson
isomorphism.
\end{proof}

\begin{remark}[Status of the analytic input]
The reduction in Theorem~\ref{thm:localmodel} is proved here in full; the only
imported statement is the equivariant $b$-Moser theorem, stated precisely above
with its hypotheses verified. We do not reprove it; it is a standard
normal-form theorem in $b$-geometry and is the $b$-analogue of the equivariant
Moser/Marle--Guillemin--Sternberg slice theorem.
\end{remark}

\section{$b$-Delzant data and the classification}

We now globalise. The subtlety flagged in the literature is that
$\overline{\mu(U)}$ is \emph{not} a bounded polytope: by
Lemma~\ref{lem:blowup} the component of $\mu$ in the modular direction diverges
to $\pm\infty$ at $Z$. The correct invariant therefore uses a
\emph{modular-compactified} moment map.

\begin{definition}[Compactified moment map]\label{def:barmu}
Fix, near each component $Z_j$ of $Z$, the normal coordinate $t_j$
of~\eqref{eq:bnf}, and choose a diffeomorphism $\tau\colon\mathbb R\to(-1,1)$,
$\tau(s)=\tfrac{2}{\pi}\arctan(s)$. Define
$\bar\mu\colon M\to\mathfrak t^*$ by setting, in the modular direction
$v^*_{Z_j}$ near $Z_j$,
\[
\langle\bar\mu,v_{Z_j}\rangle:=\lambda_j-\tau\!\bigl(-\langle\mu,v_{Z_j}
\rangle\bigr)\quad\text{(suitably oriented)},
\]
where $\lambda_j$ is the constant making $Z_j$ map to a hyperplane, and equal to
$\mu$ in the bounded directions and away from $Z$. Then $\bar\mu$ is continuous,
$T^n$-invariant, proper, sends each component $Z_j$ to (the relative interior of)
a hyperplane, and $\bar\mu(M)$ is compact.
\end{definition}

\begin{remark}
$\bar\mu$ collapses the two logarithmic ``ends'' on the two sides $t_j>0$,
$t_j<0$ of $Z_j$ onto a common finite facet. The bounded $n-1$ directions of
$\mu$ are unchanged. Thus $\bar\mu$ records exactly the bounded geometry of the
symplectic locus together with the locations of the singular hypersurfaces;
the divergent modular size is discarded (it is gauge, by translation), and the
\emph{infinitesimal} modular size $c_{Z_j}$ is retained separately. This is the
b-moment-map picture of Guillemin--Miranda--Pires--Scott.
\end{remark}

\begin{definition}\label{def:bdelzant}
A \emph{$b$-Delzant datum} in $(\mathfrak t^*,\Lambda^*)$ consists of:
\begin{enumerate}
\item a Delzant polytope $\Delta\subset\mathfrak t^*$ (convex, simple, rational,
smooth: at each vertex the $n$ primitive inward edge-vectors form a $\mathbb
Z$-basis of $\Lambda$);
\item a finite (possibly empty) set $\mathcal H$ of facets of $\Delta$,
\emph{pairwise non-adjacent} (no two facets in $\mathcal H$ share a common
$(n-2)$-face of $\Delta$), each designated \emph{singular}; and
\item for each $F\in\mathcal H$ a positive real $c_F>0$, the \emph{modular
period}.
\end{enumerate}
Two $b$-Delzant data are \emph{equivalent} if they differ by an element of
$\mathrm{AGL}(n,\mathbb Z)=\Lambda\rtimes\mathrm{GL}(\Lambda)$ acting on
$\mathfrak t^*$, carrying $\mathcal H$ to $\mathcal H$ and preserving each $c_F$.
\end{definition}

\begin{theorem}[$b$-Delzant classification; Guillemin--Miranda--Pires--Scott]
\label{thm:main}
Fix $(\mathfrak t,\Lambda)\cong(\mathbb R^n,\mathbb Z^n)$. The assignment
\[
(M,\pi,T^n,\mu)\ \longmapsto\
\bigl(\,\Delta:=\bar\mu(M),\ \mathcal H,\ (c_F)_{F\in\mathcal H}\,\bigr)
\]
(with $\bar\mu$ as in Definition~\ref{def:barmu}, $\mathcal H$ the facets of
$\Delta$ that are images of components of $Z$, and $c_F$ the modular periods of
Lemma~\ref{lem:blowup}) is a bijection between
\begin{itemize}
\item equivariant Poisson isomorphism classes of compact connected toric
$b$-symplectic $2n$-manifolds (Definition~\ref{def:toricb}), and
\item equivalence classes of $b$-Delzant data (Definition~\ref{def:bdelzant}).
\end{itemize}
\end{theorem}

\begin{proof}
We prove: well-definedness (W); that the output is a $b$-Delzant datum (D);
injectivity (I); surjectivity (S).

\medskip
\noindent\textbf{(W) Well-definedness.}
By compactness of $M$, $Z$ has finitely many components $Z_1,\dots,Z_m$, each a
compact hypersurface. By Definition~\ref{def:barmu}, $\bar\mu$ is proper,
$T^n$-invariant and continuous, so $\bar\mu(M)$ is compact. The bounded
directions of $\mu$ are unaffected, and the singular directions are collapsed to
finite facets; $\bar\mu(M)$ is independent of the auxiliary $\tau$ up to the
diffeomorphism class of its facet structure and, as a convex set (part (D)),
depends on $\mu$ only up to the translation/$\mathrm{GL}(\Lambda)$ ambiguity,
i.e.\ up to $\mathrm{AGL}(n,\mathbb Z)$. The set $\mathcal H$ and the numbers
$c_{Z_j}$ are intrinsic by Lemmas~\ref{lem:G4}--\ref{lem:blowup}. An equivariant
Poisson isomorphism $\Phi\colon M_1\to M_2$ restricts to an equivariant
symplectomorphism $U_1\to U_2$ intertwining $\mu_1,\mu_2$ up to
$\mathrm{AGL}(n,\mathbb Z)$, carries $Z^{(1)}$ to $Z^{(2)}$, and preserves the
modular periods (the modular class is a Poisson invariant, and $c_{Z_0}=r(v_{Z_0})$
is intrinsic by Lemma~\ref{lem:G4}); hence produces equivalent data.

\medskip
\noindent\textbf{(D) The output is a $b$-Delzant datum.}
We prove convexity by establishing \emph{local} convexity of $\bar\mu(M)$ from
the explicit local models (Atiyah--Guillemin--Sternberg in the interior,
Theorem~\ref{thm:localmodel} along $Z$) and then upgrading to global convexity
by the classical Tietze--Nakajima theorem; this replaces the black-box
$b$-convexity statement by the standard symplectic local theory plus a standard
convex-geometry theorem.

\emph{(D1) $\bar\mu(M)$ is compact, connected, and a finite union of local
cones.} By Definition~\ref{def:barmu} $\bar\mu$ is continuous and proper with
$M$ compact, so $\bar\mu(M)\subset\mathfrak t^*$ is compact, hence closed; it is
connected because $M$ is. We describe the germ of $\bar\mu(M)$ at each
$w\in\bar\mu(M)$.

\emph{Interior points (away from $Z$).} If $\bar\mu^{-1}(w)\subset U$, then on a
$T^n$-invariant neighbourhood $V$ of the compact fibre $\bar\mu^{-1}(w)$ the map
$\bar\mu=\mu$ is the genuine moment map of the effective Hamiltonian
$T^n$-action on the symplectic manifold $(V,\omega)$. By the local normal form
of Marle--Guillemin--Sternberg (equivariant Darboux/slice theorem), there is a
$T^n$-fixed type so that $\mu(V)$ coincides near $w$ with $w+C_w$, where $C_w$
is a rational polyhedral cone (the cone over the moment image of a linear
isotropy representation); this is the Atiyah--Guillemin--Sternberg local
convexity statement. In particular $\bar\mu(M)\cap V$ is convex near $w$ and
equals $w+C_w$ for a rational cone $C_w$.

\emph{Points over $Z$.} If $w\in\bar\mu(Z_j)=:F_j$, take a point $q\in
\bar\mu^{-1}(w)\subset Z_j$. By Theorem~\ref{thm:localmodel} a $T^n$-invariant
neighbourhood of $Z_j$ is the model $Z_j^{\mathrm{symp}}\times S^1\times
(-\varepsilon,\varepsilon)$ with $\omega_0=c_{F_j}\tfrac{dt}{t}\wedge d\phi+
\omega_{Z_j}$. In this model $\bar\mu$ factors as the product of (i) the
symplectic toric moment map $\mu_{Z_j}\colon Z_j^{\mathrm{symp}}\to F_j$ of the
quotient torus $T'=T^n/S^1_{Z_j}$ on the bounded directions $\ker v^*_{Z_j}$,
and (ii) the compactified modular coordinate
$\langle\bar\mu,v_{Z_j}\rangle=\lambda_j-\tau(-\langle\mu,v_{Z_j}\rangle)$ in the
direction $v^*_{Z_j}$, which by Lemma~\ref{lem:blowup} maps the two normal ends
$t\gtrless0$ monotonically onto a one-sided half-open interval $[\lambda_j,
\lambda_j+\delta)$ collapsing $Z_j=\{t=0\}$ to the value $\lambda_j$. Hence near
$w$,
\[
\bar\mu(M)\cap V \;=\; w+\bigl(C_w^{Z_j}\times[0,\delta)\,v^*_{Z_j}\bigr),
\]
where $C_w^{Z_j}\subset\ker v^*_{Z_j}$ is the rational cone of the symplectic
toric base $Z_j^{\mathrm{symp}}$ at the $T'$-image of $q$ (a single point of
$F_j$ if $q$ is a $T'$-fixed point, a half-space facet of $F_j$ otherwise). This
is again $w$ plus a rational polyhedral cone, and it is convex (a product of a
convex cone with a half-line). Note $F_j$ is exactly the face $\{v^*_{Z_j}=
\lambda_j\}$ of this local cone, so $F_j$ is a facet of $\bar\mu(M)$.

\emph{(D2) Global convexity (Tietze--Nakajima).} We invoke the following
classical theorem. \textbf{Tietze--Nakajima theorem} \emph{(H.\ Tietze 1928,
S.\ Nakajima 1928; see also Bj\"orck): a closed, connected subset $C$ of
$\mathbb R^N$ that is \emph{locally convex} (every point of $C$ has an open
neighbourhood $V$ with $C\cap V$ convex) is convex.} Its hypotheses hold for
$C:=\bar\mu(M)$: it is closed and connected by (D1), and locally convex by the
two cases of (D1), where each local germ $w+C_w$ resp.\ $w+(C_w^{Z_j}\times
[0,\delta)v^*_{Z_j})$ is convex. Therefore $\bar\mu(M)$ is convex. Being convex
and locally a finite rational polyhedral cone at each point (D1), with finitely
many distinct local cones (compactness, plus finiteness of fixed-point types and
of components of $Z$), $\bar\mu(M)$ is a rational convex polytope $\Delta$, and
it is simple because each local cone is simplicial (the isotropy weights at a
fixed point are linearly independent by effectiveness; the product with the
collapsed normal interval keeps simpliciality).

\emph{Smoothness at vertices (gap G3(a)).} Every vertex $w$ of $\Delta$ is the
image of a fixed point $p$ of the $T^n$-action. Such fixed points lie in $U$,
not on $Z$: indeed on $Z$ the $T^n$-action has the free circle factor
$S^1_{Z_0}$ (Lemma~\ref{lem:G4}), so no point of $Z$ is $T^n$-fixed. Hence near
$p$ the structure is genuinely symplectic and the equivariant symplectic Darboux
(Marle--Guillemin--Sternberg) theorem applies verbatim: the $n$ primitive inward
edge-vectors at $w$ are the weights of the isotropy representation at $p$ and
form a $\mathbb Z$-basis of $\Lambda$ by effectiveness. This is the Delzant
smoothness condition. (No appeal to a corner model on $Z$ is needed, precisely
because vertices avoid $Z$.)

\emph{Vertices adjacent to a singular facet.} A vertex $w$ \emph{on} a singular
facet $F=\bar\mu(Z_j)$ still has $\bar\mu^{-1}(w)$ a $T^n$-fixed point; but
fixed points avoid $Z$, while $\bar\mu^{-1}(\operatorname{relint}F)\subseteq Z_j$
and $\bar\mu^{-1}(w)$ with $w\in F$ a vertex is a fixed point of $T'$ in
$Z_j^{\mathrm{symp}}$ lifted by the collapse. Concretely, near such $w$ the
local model is $(\text{Delzant cone of }Z_j^{\mathrm{symp}}\text{ at the
}T'\text{-fixed point})\times(\text{collapsed normal})$; the
$Z_j^{\mathrm{symp}}$-cone is unimodular for the rank-$(n-1)$ lattice
$\Lambda\cap\ker v^*_{Z_j}$ by the Delzant condition in dimension $2n-2$, and
together with the primitive normal $v^*_{Z_j}$ (Lemma~\ref{lem:blowup}) the $n$
edge-vectors at $w$ form a $\mathbb Z$-basis of $\Lambda$ (a unimodular basis of
a rank-$(n-1)$ sublattice plus a vector completing it to a basis, by primitivity
of $v_{Z_j}$). Thus smoothness holds at vertices adjacent to singular facets as
well, and this uses Theorem~\ref{thm:localmodel} only in its product form,
which \emph{does} cover the corner because the corner of $\Delta$ on $F$ comes
entirely from a corner of the lower-dimensional toric base
$Z_j^{\mathrm{symp}}$.

\emph{Pairwise non-adjacency (gap G3(b)).} We must show that two distinct
components $Z_i\neq Z_j$ cannot have moment images $F_i=\bar\mu(Z_i)$,
$F_j=\bar\mu(Z_j)$ sharing a common $(n-2)$-face $G$ of $\Delta$. Suppose they
did. A point $w\in\operatorname{relint}G$ has $\bar\mu^{-1}(w)$ contained in the
closure of both $\bar\mu^{-1}(\operatorname{relint}F_i)\subseteq Z_i$ and
$\bar\mu^{-1}(\operatorname{relint}F_j)\subseteq Z_j$, and lying on the codimension-$2$
stratum of $\Delta$. The local-cone description (b-convexity) at $w$ would then
exhibit \emph{two independent} collapsed normal directions $v^*_{Z_i},v^*_{Z_j}$,
i.e.\ two transverse singular hypersurfaces meeting over $\bar\mu^{-1}(w)$. But
the preimage $\bar\mu^{-1}(\operatorname{relint}G)$ is a $T^n$-invariant subset
on which $\pi$ would have corank $\ge2$: along it $\pi^n$ vanishes ``from two
sides,'' contradicting that $\pi^n$ vanishes \emph{transversally} on the smooth
hypersurface $Z$ (Definition~\ref{def:bpoisson}), whose components $Z_i,Z_j$
are by definition disjoint and each of corank exactly $1$. Concretely: at a
point $q$ in the alleged common stratum, two distinct local defining functions
$t_i,t_j$ would both vanish with linearly independent differentials and $\pi$
would degenerate in both normal directions, forcing $\operatorname{corank}
\pi(q)\ge2$, impossible for a $b$-Poisson structure. Hence $F_i,F_j$ share no
$(n-2)$-face; the singular facets are pairwise non-adjacent. (This corrects the
earlier ``intersection of $Z_i,Z_j$'' phrasing: the $Z$'s are disjoint, and the
real obstruction is to their \emph{images} being adjacent, which is what the
corank argument above rules out.)

Thus $(\Delta,\mathcal H,(c_F))$ is a $b$-Delzant datum.

\medskip
\noindent\textbf{(I) Injectivity.}
Suppose $(M_1,\pi_1)$, $(M_2,\pi_2)$ give equivalent data; after acting by
$\mathrm{AGL}(n,\mathbb Z)$ assume identical $(\Delta,\mathcal H,(c_F))$.

\emph{Interior.} The complement of the singular facets, $\Delta^\circ:=
\Delta\setminus\bigcup_{F\in\mathcal H}F$, is a (non-compact, ``cylindrical
ended'') Delzant region. On $U_i=\bar\mu_i^{-1}(\Delta^\circ\cup\operatorname
{int}\Delta)$ the action is Hamiltonian toric with the same moment image. By the
uniqueness half of Delzant's theorem in the non-compact toric case (Karshon--Lerman,
``Non-compact symplectic toric manifolds,'' SIGMA 2015: a symplectic toric
manifold, compact or not, is determined up to equivariant symplectomorphism by
its moment image as a manifold-with-corners with its characteristic lattice
data), there is an equivariant symplectomorphism $\Psi\colon U_1\to U_2$ with
$\mu_2\circ\Psi=\mu_1$. The hypotheses of Karshon--Lerman hold because, by
Lemma~\ref{lem:blowup} and Theorem~\ref{thm:localmodel}, each $U_i$ is a
symplectic toric manifold whose moment map is proper onto $\Delta^\circ$ with
the explicit logarithmic cylindrical end normal to each $F$, so the
characteristic data and the ``completeness at the ends'' required by their
theorem are present and identical for $i=1,2$.

\emph{Across $Z$ and elimination of gluing moduli (gap G2).} Near each
$Z_j^{(i)}$ Theorem~\ref{thm:localmodel} identifies a neighbourhood with the
same model $\bigl(Z_j^{\mathrm{symp}}\times S^1\times(-\varepsilon,\varepsilon),
\,c_{F_j}\tfrac{dt}{t}\wedge d\phi+\omega_{Z_j}\bigr)$, because: (a) the base
$Z_j^{\mathrm{symp}}$ is the symplectic toric $(2n-2)$-manifold with Delzant
polytope $F_j$, hence \emph{determined} (Delzant in dimension $2n-2$) and equal
for $i=1,2$; (b) the modular period $c_{F_j}$ agrees; and (c) the principal
$S^1_{Z_j}$-bundle structure is the \emph{trivial} bundle in the model and the
model fixes the connection $\theta=c_{F_j}d\phi$. We must check there is no
residual gluing modulus. The only possible moduli when assembling $M$ from the
two ``ends'' $\{t>0\}$, $\{t<0\}$ across $Z_j$ are (i) a deck/automorphism of
the symplectic toric base $Z_j^{\mathrm{symp}}$ compatible with the
$T'$-moment data, and (ii) a $b$-cohomology class measuring the difference of
the two $b$-symplectic structures. For (i): an automorphism of a symplectic
toric manifold fixing its moment map is the identity (Delzant uniqueness with
fixed moment map), so the only base automorphism compatible with $F_j$ is the
identity; the modular $S^1$-direction is rigid because $\theta$ is fixed to
$c_{F_j}d\phi$. For (ii): the difference $\alpha:=\omega_2-\omega_1$ of the two
$b$-symplectic forms, after applying $\Psi$ on the overlap, is a $T^n$-invariant
closed $b$-$2$-form on the model collar
$N_j:=Z_j^{\mathrm{symp}}\times S^1\times(-\varepsilon,\varepsilon)$. By the
Mazzeo--Melrose decomposition, the $b$-de~Rham cohomology of $N_j$ splits as
$H^k_b(N_j)\cong H^k(N_j)\oplus H^{k-1}(Z_j)$, where the second summand is the
\emph{residue}; the class of a closed $b$-$2$-form is determined by its ordinary
class in $H^2(N_j)$ together with its residue $r(\alpha)\in H^1(Z_j)$. Here the
residue of $\omega_i$ along $Z_j$ is $[c_{F_j}d\phi]\in H^1(Z_j)$ for both $i$
(Lemma~\ref{lem:blowup} with the same period $c_{F_j}$), so $r(\alpha)=0$;
moreover $N_j$ deformation retracts onto $Z_j$ and $\omega_1,\omega_2$ restrict
to the same leafwise/base symplectic data $\omega_{Z_j}$ there, so the ordinary
class $[\alpha]\in H^2(N_j)\cong H^2(Z_j)$ also vanishes. Thus $\alpha$ is
$b$-exact: $\alpha=d_b\eta$ for some $b$-$1$-form $\eta$ on $N_j$. Averaging
$\eta$ over the compact group $T^n$ (using $T^n$-invariance of $\alpha$ and
$d_b$-equivariance) we may take $\eta$ to be $T^n$-invariant, and subtracting a
$d_b$-closed correction we may take $\eta$ to vanish on $Z_j$. The linear path
$\omega_s:=\omega_1+s\,\alpha$ ($s\in[0,1]$) consists of $b$-symplectic forms
(nondegeneracy is open and both endpoints share the same residue and leafwise
data, so $\omega_s$ is nondegenerate on the collar after shrinking
$\varepsilon$), and the equivariant $b$-Moser theorem
(Theorem~\ref{thm:localmodel}) applied to $\eta$ — whose Moser vector field
$Y_s$, defined by $\iota_{Y_s}\omega_s=-\eta$, is $T^n$-invariant, tangent to
$Z_j$, and complete because $Z_j$ is compact and $\eta$ vanishes on $Z_j$ to the
order making the flow well-defined across the pole — yields a $T^n$-equivariant
isotopy fixing $Z_j$ and carrying $\omega_1$ to $\omega_2$. There is therefore no
free parameter in the gluing. Patching $\Psi$ (on the interior) to these
boundary equivariant Poisson isomorphisms by a $T^n$-invariant partition of
unity subordinate to $\{U_i$-interior$\}\cup\{$model collars$\}$, and removing
the resulting discrepancy by one final equivariant Moser isotopy on the compact
$M$, yields a global equivariant Poisson isomorphism $M_1\cong M_2$. Here the
discrepancy $\omega'-\omega$ is a $T^n$-invariant closed $b$-$2$-form supported
in the symplectic region away from $Z$ (the boundary maps were already exact
matches near $Z$ by the previous paragraph), with the same cohomology class
because the two structures induce the same moment data on $\Delta$; hence it is
$d_b$ of a $T^n$-invariant $1$-form (average a primitive, which exists by the
ordinary Poincar\'e lemma on the contractible-fibre region and the vanishing of
the residue away from $Z$), and the equivariant Moser argument on the compact
symplectic locus (standard, since here there is no pole) completes the global
isomorphism.

\medskip
\noindent\textbf{(S) Surjectivity (construction).}
Given $(\Delta,\mathcal H,(c_F))$, build $M$ explicitly.

\emph{Step 1 (Delzant in dimension $2n$).} Realise $\Delta$ by its facet
inequalities $\langle\xi,a_i\rangle\ge\lambda_i$, $i=1,\dots,d$, with
$a_i\in\Lambda$ primitive inward normals. Let $\beta\colon\mathbb R^d\to
\mathfrak t$, $e_i\mapsto a_i$; let $N=\ker(\exp\colon T^d\to T^n)$ act on
$\mathbb C^d$ with standard moment map. The reduction
$X_\Delta=\mu_{\mathbb C^d}^{-1}(\lambda)/N$ is the compact symplectic toric
manifold with polytope $\Delta$ and residual action of $T^n=T^d/N$ with moment
map $\mu_\Delta$ onto $\Delta$.

\emph{Step 2 ($b$-modification along each singular facet).} Fix $F\in\mathcal H$
with primitive normal $a_F$ and period $c_F$. Let $H_F:=\mu_\Delta^{-1}(F)$, a
$T^n$-invariant codimension-$1$ submanifold (the preimage of a facet), and let
$s$ be the smooth moment coordinate $s=\langle\mu_\Delta,v_F\rangle-\lambda_F$
normal to $F$, where $v_F\in\Lambda$ is primitive dual to $a_F$; thus $H_F=
\{s=0\}$ and $X_\Delta$ is symplectic with $\omega_\Delta=ds\wedge d\phi_F+
\omega_F^{\mathrm{base}}+\cdots$ in a collar, $\phi_F$ the angle of
$\exp(\mathbb R v_F)$ and $\omega_F^{\mathrm{base}}$ the symplectic form of the
toric facet manifold $X_F$ (the symplectic toric $(2n-2)$-manifold with polytope
$F$, by Delzant in dimension $2n-2$). On a collar
$U_\epsilon=\{0<|s|<\epsilon\}$ replace the symplectic bivector by the
$b$-bivector obtained from
\[
\omega_F^{b}:=c_F\,\frac{dt}{t}\wedge d\phi_F+\omega_F^{\mathrm{base}},\qquad
t:=\exp(s/c_F),
\]
so that $\tfrac{dt}{t}=\tfrac1{c_F}ds$ and $c_F\tfrac{dt}{t}\wedge d\phi_F=
ds\wedge d\phi_F$: the two forms \emph{agree} on the overlap $0<|s|<\epsilon$,
so the modification is the identity there and the global bivector is smooth.
After the substitution, $t$ ranges over $(0,\infty)$ as $s\in(-\infty,\infty)$
in the model; we glue the two collars $s>0$ and $s<0$ across the new
hypersurface $Z_F:=\{t=0\}$ using the local model of
Theorem~\ref{thm:localmodel} with base $X_F$ and period $c_F$. Pairwise
non-adjacency of the facets in $\mathcal H$ guarantees the submanifolds $H_F$
are pairwise disjoint (their facets share no $(n-2)$-face, so their collars are
disjoint), so the modifications are independent.

\emph{Verification of the structure (gap G6).} (a) \emph{$b$-Poisson:} in the
collar, $\omega_F^b$ has $b$-symplectic dual $\pi_F$ with $\pi_F^n$ vanishing
exactly to first order along $\{t=0\}$ (the factor $t$ in $t\,\partial_t\wedge
\partial_{\phi_F}$ has transverse zero), so $\pi$ is $b$-Poisson with $Z=
\bigsqcup_F Z_F$. (b) \emph{Smoothness/matching:} by the chosen substitution
$t=\exp(s/c_F)$ the modified form equals $\omega_\Delta$ on $\{|s|>0\}$ and the
gluing uses the exact product model, so $(M,\pi)$ is a smooth manifold and
$\pi$ a smooth bivector. (c) \emph{Effective Poisson $T^n$-action:} the
modification is $T^n$-equivariant (it only reparametrises the $s$-direction,
which is the moment direction of the subcircle $\exp(\mathbb R v_F)$, and
preserves $\omega_F^{\mathrm{base}}$), and effectiveness is inherited from
$X_\Delta$. (d) \emph{Compact connected:} $M$ is obtained from compact
$X_\Delta$ by an operation supported on collars and replacing each $H_F$ by the
hypersurface $Z_F$ (a finite, codimension-$1$ surgery that preserves
compactness and connectedness, since $\mathcal H$ is finite and each
modification is a homeomorphism onto its image away from $Z_F$). Recovering the
datum: by construction $\bar\mu(M)=\Delta$, the components of $Z$ are the
$Z_F$ with $\bar\mu(Z_F)=F$, and the modular periods are the prescribed $c_F$
(by Lemma~\ref{lem:blowup} applied to $\omega_F^b$). When $\mathcal H=
\varnothing$ this is exactly Delzant's construction.

\medskip
(W),(D),(I),(S) together establish the bijection.
\end{proof}

\section{Answer to the question}

\begin{corollary}
The Delzant classification \emph{does} extend to the Poisson setting, but only
after restricting to a class of Poisson structures whose singular locus is
controlled. The right class is the $b$-symplectic one: compact connected toric
$b$-symplectic $2n$-manifolds are in bijection with $b$-Delzant data --- a
Delzant polytope $\Delta$, a set $\mathcal H$ of pairwise non-adjacent facets
marked as singular, and a positive modular period $c_F$ for each $F\in\mathcal
H$. The convex-geometric object is the compactified moment image $\Delta=
\bar\mu(M)$ together with these decorations; the new invariants relative to the
symplectic case are exactly $\mathcal H$ and $(c_F)$, encoding where and how
strongly $\pi$ degenerates. For Poisson structures whose degeneracy is \emph{not}
of $b$-type, Proposition~\ref{prop:noinvariant} shows no such clean
classification can exist.
\end{corollary}

\begin{remark}[Scope and honest limitations]
The theorem is the Delzant-type classification of
Guillemin--Miranda--Pires--Scott. We have derived in full the structural
lemmas (Lemmas~\ref{lem:Z}, \ref{lem:G4}, \ref{lem:blowup}), the reduction in
the local normal form (Theorem~\ref{thm:localmodel}), the well-definedness,
the smoothness and non-adjacency in part (D), the elimination of gluing moduli
in part (I), and the explicit construction in part (S). The convexity step (part (D)) is reduced here to the explicit local models
(Atiyah--Guillemin--Sternberg in the interior, Theorem~\ref{thm:localmodel}
along $Z$) together with the classical Tietze--Nakajima convexity theorem, so it
is not black-boxed. The one genuinely analytic input we \emph{cite, not reprove}
is the equivariant $b$-Moser normal-form theorem (used in
Theorem~\ref{thm:localmodel} and in part (I)); we stated it precisely and
verified its hypotheses (compact invariant $Z$, equal residues, matching
restriction data, completeness of the Moser field) in our setting. For higher-order vanishing of $\pi^n$
($b^k$-structures, $k\ge2$) the single period $c_F$ is replaced by a $k$-jet of
modular data, and Proposition~\ref{prop:noinvariant}(2) shows that for $k\ge2$
the moment image alone is no longer a complete invariant; the $b^k$-analogue of
the present theorem requires the $b^k$-Moser theorem and is not treated here.
\end{remark}

\end{document}
